\documentclass[xcolor={dvipsnames,table}]{beamer}
\mode<presentation>{\usetheme{Warsaw}\usecolortheme{crane}}
\usepackage{centernot}
\usepackage{graphicx}
\usepackage{geometry}
\usepackage{tikz}
\usetikzlibrary{shadows}

\usefonttheme{serif}
\usepackage[utf8]{inputenc}
\usepackage[english]{babel}
\usepackage{lmodern}
\usepackage[T1]{fontenc}
\usepackage[babel=true]{microtype}

\title{UNIX Weapons School---Computer Architecture}
\date{}
\author{Nick Black for the\\
Georgia Institute of Technology
}

\begin{document}

\begin{frame}
\titlepage
\begin{center}
\includegraphics[scale=0.33]{images/uws.png}\\
\vspace{.1in}
\tiny{copyright \copyright\ 2013}\\
\includegraphics[scale=.25]{images/cc-logo.pdf}
\includegraphics[scale=.25]{images/cc-new.pdf}
\includegraphics[scale=.25]{images/cc-share.pdf}\\
\tiny{creative commons 3.0 share-alike attribution license}
\end{center}
\end{frame}

\begin{frame}{Recommended Reading}
\small
\begin{itemize}
\item Andi Kleen. ``Linux multi-core scalability.'' Linux Kongress 2009.
\item Ulrich Drepper. ``What Every Programmer Needs to Know About Memory.'' In 8 parts. Linux Weekly News (2007).
\item Paul McKenney. ``Transactional Memory Everywhere.''\\
http://paulmck.livejournal.com (2012).
\item John Shen and Mikko Lipasti. \textit{Modern Processor Design} (2004).
\item \textit{Intel 64 and IA-32 Architectures Optimization Manual}. Intel.
\item Agner Fog. \textit{The microarchitecture of Intel, AMD and VIA CPUs: An optimization guide for assembly programmers and compiler makers.}
Regularly updated. http://www.agner.org.
\end{itemize}
\end{frame}

\begin{frame}{The purpose of a systems programmer}
\textbf{Given system $S$, and problem $P$,\\
we will implement $P$,\\
using the resources of $S$,\\
subject to some constraint.}\\
\vspace{.25in}
That constraint is typically to minimize either:
\begin{itemize}
\vspace{.05in}
\item time to completion,
\vspace{.05in}
\item power to completion, or
\vspace{.05in}
\item $\frac{\text{time to completion}}{\text{utilization}}$
(subject to some minimum utilization).
\end{itemize}
\vspace{.25in}
Proving that these constraints have been met is clearly
dependent upon the details of both $P$ and $S$.
\end{frame}

\begin{frame}{Abstraction}
\begin{itemize}
\item ``The purpose of abstraction is not to be vague, but to create a new semantic level in which one can be absolutely precise.''\\
\hfill--Edsger Dijkstra
\vfill
\item ``First you learn the value of abstraction. Then you learn the cost of abstraction. Then you're ready to engineer.''\\
\hfill--Kent Beck
\vfill
\item We will draw the line at the electric level.
\end{itemize}
\end{frame}

\begin{frame}{Complementary metal-oxide semiconductors I}
Draw \textit{greatest} power only when changing state\ldots
\begin{equation}
P_{dynamic} = \alpha C_LV_{DD}^{2}\nu
\end{equation}
\begin{description}
\item[$P$] Dynamic power ($W = J/s = \frac{m^{2}\cdot kg}{s^{3}}$)
\item[$C_L$] Load capacitances ($F = \frac{A^{2}\cdot s^{4}}{m^{2}\cdot kg}$)
\item[$\alpha$] Activity factor
\item[$V_{DD}$] Supply voltage ($V = \frac{m^{2}\cdot kg}{A\cdot s^{3}}$)
\item[$\nu$] Frequency ($1/s$)
\end{description}
\vfill
\textbf{NB:} $P_{dynamic}$ increases with the \textit{square} of $V_{DD}$.
\end{frame}

\begin{frame}{Complementary metal-oxide semiconductors II}
\ldots but always consume \textit{some} power:
\begin{equation}
P_{static} = I_{static}\cdot V_{DD}
\end{equation}
\begin{description}
\item[$P$] Static power ($W = V\cdot A$)
\item[$I_{static}$] Static current ($I_{subth} + I_{tunnel} + I_{leakage}, A$)
\item[$V_{DD}$] Supply voltage ($V$)
\end{description}
\vfill
\textbf{Implication:} power draw is independent of actual state.
\end{frame}

\begin{frame}{Instructions I}
A system starts in some initial state, and while undisturbed, evolves
in discrete time according to some closed primitive recursive function on that state (aka Church-Turing).\\
\vspace{.25in}
Each time step is a \textit{cycle}. 1GHz$\implies$ 1ns cycle.\\
\vspace{.25in}
``Running a program'' is the act of denoting one or more words of this state
as ``instructions'', causing them to drive control flow.
\vspace{.25in}
\begin{itemize}
\item Static instructions: Instructions in a sequence
\item Dynamic instructions: Instructions \textit{executed} in a sequence
\item IPC: Instructions per cycle
\item CPI: Cycles per instruction
\end{itemize}
\end{frame}

\begin{frame}{Instructions II}
Instructions lie at the boundary between those abstractions we \textit{control} and those
we merely \textit{exploit}.
\begin{equation}
P_{time} = Inst_{dynamic}\cdot CPI_{avg}\cdot\nu
\end{equation}
The means of accelerating our program are finite:
\begin{enumerate}
\item Reduce cycle time (better $\mu$architecture, better materials)
\item Require fewer instructions (better code, better algorithms, more powerful instructions, \textbf{BQP})
\item Reduce cycles per instruction (better $\mu$architecture, better code)
\item Add time (increased high-level parallelism)
\item Dilate time (relativistic motion)
\end{enumerate}
\end{frame}

\begin{frame}{Special registers}
Often read-only/privileged/require special instructions.
\tiny{
\begin{itemize}
\item $EIP$: Read-only. Instruction pointer
\item $EFLAGS$: Flags register. Contains IOPL/CPL (\texttt{POPF}/\texttt{IRET})
\item $CR_0$: Machine Status Control Word (\texttt{LMSW}, 286+)
\item $CR_2$: Read-only. Page fault addresses (386+)
\item $CR_3 (PDBR)$: Page Directory Base Reg. (386+)
\item $MXCSR$ SSE Control Status Reg. (\texttt{LD/STMXCSR}, PIII+)
\item $DR_0, DR_1, DR_2, DR_3, DR_7$: Debug Reg.
\item $CS, DS, ES, FS, GS, SS$: Segment Reg.
\item $TR$: Task Register (\texttt{LTR}/\texttt{STR})
\item $GDTR$: Global Descriptor Table Reg. (\texttt{LGDT}/\texttt{SGDT}, 286+)
\item $LDTR$: Local Descriptor Table Reg. (\texttt{LLDT}, 286+)
\item $IDTR$: Interrupt Descriptor Table Reg. (\texttt{LIDT}/\texttt{SIDT}, 286+)
\item $TPR$: Task Priority Reg.
\item $PPR$: Process Priority Reg.
\item $TSC$: Timestamp Counter (\texttt{RDTSC}, Pentium+)
\item $PMC_n$: Performance Monitoring Counters (\texttt{RDPMC}, MMX+)
\item $MSR_n$: Model-Specific Reg. (\texttt{RD/WRMSR}, Pentium+)
\end{itemize}
}
\end{frame}

\begin{frame}{General-purpose registers}
\begin{itemize}
\item Read ports pace superscalability
\item Write ports pace instruction retire
\item Register width paces bit parallelism
\item Architectural registers pace memory accesses
\item Physical registers pace OOO hiding of antidependency
% particularly useful for EFLAGS!
\end{itemize}
\end{frame}

\begin{frame}{Register file (indexed access)}
\begin{itemize}
\item Architectural registers are exposed as PRAM or a stack
\item SRAM cells + read/write lines + decoder tree + sense amps
\item One internal bit line per bit of read port
\item Two internal bit lines per bit of write port
\item One word line per entry
\item Transistor area grows linearly with number of ports
\item Wire pitch area grows with square of number of ports
\end{itemize}
\end{frame}

\begin{frame}{Fundamental theorem of memory starvation}
We can only hit peak during sequences of code which fully
utilize arithmetic resources. Registers must provide throughput
at least equivalent to the processor's arithmetic throughput.\footnote{\tiny{In a load/store (classic RISC) architecture, arithmetic instructions operate
only on registers. In a machine admitting memory addresses as inputs to
arithmetic operands, such instructions will exhibit greater latency and lesser throughput
than their register-direct counterparts.}}\\
\vspace{.25in}
\textbf{THUS,} a memory lacking throughput greater than or equal to the ratio of
arithmetic/register throughput to arithmetic intensity (cycles per word), or having latency
greater than or equal to arithmetic latency (seconds per word), will not be able
to sustain peak.
\end{frame}

\end{document}
